% Introduction
Many Internet of Things applications depend on data that is not only correct but
also fresh. In precision agriculture, environmental monitoring, disaster
response, and industrial inspection, a control loop or an operator acts on the
most recent sensor reading, so a stale reading can be worse than no reading at
all. Ground sensors in these settings are often sparse, spread over a wide area,
and far from any gateway, which makes continuous fresh collection difficult with
fixed infrastructure alone.

Unmanned aerial vehicles (UAVs) are attractive data collectors for such
deployments: they fly close to ground sensors, establish strong line-of-sight
links, and reach areas without fixed infrastructure. When the collected data is
time-sensitive, the relevant metric is the Age of Information (AoI)
\cite{kaul2012aoi}, that is, the time elapsed since the freshest delivered
update was generated. Keeping the network AoI low requires a UAV to revisit each
sensor often, which conflicts with the limited battery of a rotary-wing UAV: the
platform must periodically leave the field to recharge at a charging station.

A single UAV can only keep a small area fresh, so practical deployments use a
\emph{swarm}. A swarm sharing a small number of charging stations creates a
problem that a single UAV never faces: the stations become a \emph{contended}
resource. When several UAVs need to recharge at the same time, they compete for
a limited number of charging ports and must wait, and this waiting time inflates
the AoI. The system designer must therefore decide jointly where to place the
charging stations, how many ports to install at each, how the swarm should fly,
and which UAV recharges where.

\subsection{Limitations of prior work}
Existing studies address parts of this problem but not the whole. The closest
work jointly optimizes charging-station placement and trajectory, but for a
\emph{single} UAV, so there is no contention \cite{base2024}. Multi-UAV studies
that model charging either assume a single fixed facility \cite{wei2024}, use
mid-flight laser charging with fixed emitters and no queue \cite{sun2025laser},
or use mobile ground vehicles that follow the UAVs and thereby remove the
fixed-station contention \cite{oubbati2024ugv,zhao2025aoi}. Swarm-AoI studies
that ignore energy altogether \cite{eldeeb2023swarm} cannot capture recharging
at all. A separate line models UAV charging stations as a queue with waiting
time \cite{ev2022sharing}, but without any AoI objective, trajectory, or
placement. Crucially, a naive queueing model here is misleading: the charging
system is a \emph{closed} population of only $M$ UAVs, so an open $M/M/c$ queue
substantially overpredicts the wait, as we show against a discrete-event
simulation.

\subsection{Contributions}
\begin{itemize}
\item We formulate the first joint charging-station \emph{placement}, port
\emph{capacity}, swarm \emph{trajectory}, and \emph{assignment} problem for
peak-AoI fresh data collection with a contended, multi-server charging resource,
as a mixed-integer nonlinear program (Section~\ref{sec:model}).
\item We model the charging stations as a \emph{finite-source} (machine-repair)
queue and validate it against a discrete-event simulation. The finite-source
model matches the simulation within about $4\%$ and is a conservative upper
bound on the real wait, whereas the open $M/M/c$ overpredicts by about $6\times$
and distorts the design decisions.
\item We prove two structural results
(Section~\ref{sec:analysis}): a \emph{congestion threshold} that makes the AoI
non-monotone in the swarm size, with a provisioning law linking the optimal
swarm size to the charging capacity, and a \emph{placement inversion} showing
that the coverage-optimal placement is AoI-suboptimal under contention.
\item We design a block-coordinate solver with a GPU close-enough-TSP trajectory
optimizer (Section~\ref{sec:solver}) and evaluate it against single-UAV,
single-station, coverage-optimal, and round-robin baselines
(Section~\ref{sec:exp}). The proposed method reduces peak AoI significantly, and
we report the results honestly, including the regime where a single pooled
station is already strong. We further show that the greedy solver is within
$0.1\%$ of the exact optimum on small instances and that its advantage is robust
to the collection radius, the charge time, and non-uniform sensor distributions.
\end{itemize}

The remainder of the paper is organized as follows.
Section~\ref{sec:related} reviews related work,
Section~\ref{sec:model} presents the model and problem,
Section~\ref{sec:analysis} the structural analysis,
Section~\ref{sec:solver} the solver, and
Section~\ref{sec:exp} the experiments. Section~\ref{sec:concl} concludes.
